\chapter{Gaps above the kernel and the martingale method}
\label{ch:martingale_gaps}

This chapter collects lower bounds for the gap of a positive operator above its
kernel on a finite-dimensional complex Hilbert space, and Nachtergaele's martingale
method~\cite{Nachtergaele1996Spectral}, which bounds the gap of a frustration-free
Hamiltonian from below by local gaps and an overlap estimate for nested ground
spaces. It closes with the contraction of a sweep of projections~\cite{Anshu2016Simple}
and with commuting parent Hamiltonians on three regions of a system of
sites~\cite[Appendix~D.2]{Cirac2016MPDO_arXiv}, in the notation of
Chapters~\ref{ch:sites}, \ref{ch:dynamics} and~\ref{ch:hamiltonians}.

An operator is a linear map. An operator $H$ is positive, $H\ge0$, when it is
symmetric and $\operatorname{Re}\langle Hv,v\rangle\ge0$ for every $v$; $H\le K$ means
$K-H\ge0$; an orthogonal projection is a symmetric idempotent. For $\gamma\in\R$, the
operator $H$ has norm gap $\gamma$ when $\gamma\lVert v\rVert\le\lVert Hv\rVert$ for every
$v\perp\ker H$. For a positive semidefinite matrix this is the gap at energy $0$ of
Theorem~\ref{thm:hamiltonians_kernel_gap}.

\section{Gaps of positive operators above the kernel}
\label{sec:martingale_gaps_kernel}

\subsection{Quadratic forms and norms}
\label{sec:martingale_gaps_forms}

\begin{lemma}[A quadratic-form inequality gives a norm gap]
    \label{lem:gap_quadratic_form_to_norm}
    \lean{FrustrationFree.spectralGap_of_martingale_of_finiteDimensional}
    \leanok
    Let $H\ge0$ on a finite-dimensional complex Hilbert space, and let $\gamma>0$
    satisfy $\gamma\operatorname{Re}\langle Hv,v\rangle\le\lVert Hv\rVert^2$ for every
    $v$, that is, $H^2\ge\gamma H$. Then $H$ has norm gap $\gamma$. In particular every
    positive eigenvalue of $H$ is at least $\gamma$. Positivity is needed: $H=-1$ with
    $\gamma=2$ satisfies the inequality but not the conclusion.
\end{lemma}
\begin{proof}\leanok
    Let $H\psi_k=\lambda_k\psi_k$ for an orthonormal eigenbasis. The hypothesis at
    $\psi_k$ gives $\gamma\lambda_k\le\lambda_k^2$, so $\lambda_k\le0$ or
    $\lambda_k\ge\gamma$, and $H\ge0$ excludes $\lambda_k<0$. Every $v\perp\ker H$ is
    $v=\sum_kc_k\psi_k$ with $c_k=0$ whenever $\lambda_k=0$, so
    $\lVert Hv\rVert^2=\sum_k\lvert c_k\rvert^2\lambda_k^2\ge\gamma^2\lVert v\rVert^2$.
\end{proof}

\begin{lemma}[A norm gap gives the quadratic-form inequality]
    \label{lem:gap_norm_to_quadratic_form}
    \lean{LinearMap.IsPositive.quadraticForm_sq_ge_of_norm_gap}
    \leanok
    Let $H\ge0$ on a finite-dimensional complex Hilbert space, and let $\gamma\ge0$
    be a norm gap of $H$. Then
    $\gamma\operatorname{Re}\langle Hv,v\rangle\le\operatorname{Re}\langle Hv,Hv\rangle$
    for every $v$.
\end{lemma}
\begin{proof}\leanok
    Put $w=v-P_{\ker H}v$. Then $w\perp\ker H$ and $Hw=Hv$, and symmetry gives
    $\langle Hv,v\rangle=\langle Hv,w\rangle$. Hence
    $\gamma\operatorname{Re}\langle Hv,v\rangle\le\gamma\lVert Hv\rVert\lVert w\rVert
        \le\lVert Hv\rVert\lVert Hw\rVert=\operatorname{Re}\langle Hv,Hv\rangle$.
\end{proof}

\begin{lemma}[An energy bound gives a norm bound]
    \label{lem:gap_energy_to_norm}
    \lean{FrustrationFree.norm_lower_bound_of_energy_lower_bound}
    \leanok
    Let $H$ be an operator on a complex inner product space. If
    $\kappa\lVert v\rVert^2\le\operatorname{Re}\langle Hv,v\rangle$ for every
    $v\perp\ker H$, then $H$ has norm gap $\kappa$. Neither finite dimension nor
    positivity is needed.
\end{lemma}
\begin{proof}\leanok
    By the Cauchy--Schwarz inequality,
    $\kappa\lVert v\rVert^2\le\operatorname{Re}\langle Hv,v\rangle
        \le\lVert Hv\rVert\lVert v\rVert$. This is the claim for $v=0$, and gives it
    after division by $\lVert v\rVert$ when $v\neq0$.
\end{proof}

\begin{lemma}[A linear map is bounded below off its kernel]
    \label{lem:gap_orthogonal_kernel}
    \lean{LinearMap.exists_pos_mul_norm_le_of_mem_orthogonal_ker}
    \leanok
    Let $f$ be a linear map from a finite-dimensional real or complex inner product
    space to a normed space. There is $\delta>0$ with
    $\delta\lVert v\rVert\le\lVert fv\rVert$ for every $v\perp\ker f$.
\end{lemma}
\begin{proof}\leanok
    The restriction of $f$ to $(\ker f)^\perp$ is injective, and an injective linear
    map on a finite-dimensional space is bounded below: it is antilipschitz with
    some constant $K>0$, and $\delta=K^{-1}$ works.
\end{proof}

\begin{lemma}[One constant for all indices]
    \label{lem:gap_eventual_constant}
    \lean{Nat.exists_pos_forall_of_eventually}
    \leanok
    Let $P(N,\gamma)$ be a property of $N\in\N$ and $\gamma\in\R$ that is preserved
    by decreasing $\gamma$. If every $N$ has some $\eta_N>0$ with $P(N,\eta_N)$, and
    one $\delta>0$ satisfies $P(N,\delta)$ for all $N\ge M$, then one $\gamma>0$
    satisfies $P(N,\gamma)$ for every $N$.
\end{lemma}
\begin{proof}\leanok
    Take $\gamma=\min\{\delta,\eta_0,\dots,\eta_{M-1}\}$, the minimum of finitely many
    positive numbers.
\end{proof}

\begin{lemma}[Kernels and gaps under a two-sided comparison]
    \label{lem:gap_positive_comparison}
    \lean{LinearMap.IsPositive.ker_le_of_le,
        LinearMap.IsPositive.ker_eq_of_smul_le_of_le_smul,
        LinearMap.IsPositive.norm_gap_of_smul_le_of_le_smul}
    \leanok
    If $0\le P\le Q$, then $\ker Q\subseteq\ker P$. Consequently, if $P,Q\ge0$ and
    $a,b,c>0$ satisfy $aP\le bQ\le cP$, then $\ker P=\ker Q$. On a finite-dimensional
    space, a norm gap $\gamma>0$ of $P$ then gives the norm gap $a\gamma/b$ of $Q$.
\end{lemma}
\begin{proof}\leanok
    The quadratic form of a positive operator vanishes only on its kernel; evaluating
    the order relations at kernel vectors gives the inclusions. On the common
    orthogonal complement of the kernel, the norm gap of $P$ gives
    $\operatorname{Re}\langle Pv,v\rangle\ge\gamma\lVert v\rVert^2$, hence
    $b\operatorname{Re}\langle Qv,v\rangle\ge a\gamma\lVert v\rVert^2$, and
    Lemma~\ref{lem:gap_energy_to_norm} gives the norm gap of $Q$.
\end{proof}

\subsection{Reducing projections}
\label{sec:martingale_gaps_reducing}

\begin{lemma}[Commuting projections]
    \label{lem:gap_commuting_projections}
    \lean{Submodule.commute_starProjection_of_invariant,
        LinearMap.IsSymmetricProjection.mul_of_commute,
        LinearMap.IsSymmetricProjection.one_sub_mul_le_add_one_sub,
        LinearMap.IsSymmetricProjection.one_sub_mul_mul_le}
    \leanok
    A symmetric operator $T$ with $TS\subseteq S$ commutes with the orthogonal
    projection $P_S$. If $P$ and $Q$ are commuting orthogonal projections, then $PQ$
    is an orthogonal projection and $1-PQ\le(1-P)+(1-Q)$. If $P,R,C$ are pairwise
    commuting orthogonal projections, then $1-RCP\le(1-P)+(1-R)+(1-C)$.
\end{lemma}
\begin{proof}\leanok
    The identity $\langle Tu,w\rangle=\langle u,Tw\rangle$ makes $S^\perp$ invariant
    under $T$ as well, so $TP_S=P_ST$. For commuting $P,Q$ the product $PQ$ is
    symmetric and idempotent, and
    $(1-P)+(1-Q)-(1-PQ)=(1-P)(1-Q)\ge0$ as a product of commuting orthogonal
    projections. For three projections, $RC$ is an orthogonal projection commuting
    with $P$, so $1-RCP\le(1-RC)+(1-P)\le(1-R)+(1-C)+(1-P)$.
\end{proof}

\begin{lemma}[The kernel from a reducing projection]
    \label{lem:gap_reducing_projection_kernel}
    \lean{LinearMap.apply_eq_of_le_of_sub_le_of_apply_eq_zero,
        LinearMap.IsSymmetricProjection.apply_eq_self_of_one_sub_le_twice_of_mem_ker,
        LinearMap.IsPositive.ker_eq_of_le_of_eq_on_projection}
    \leanok
    \uses{lem:gap_positive_comparison}
    Let $H,K,T,Q$ be operators on a complex inner product space.
    \begin{enumerate}
        \item If $H\le K$, $K-H\le T$ and $Tv=0$, then $Hv=Kv$.
        \item If $Q$ is an orthogonal projection, $1-Q\le2H$ and $Hv=0$, then $Qv=v$.
        \item If $H\ge0$, $Q$ is an orthogonal projection, $H\le K$, $HQ=KQ$ and
              $1-Q\le2H$, then $\ker H=\ker K$.
    \end{enumerate}
\end{lemma}
\begin{proof}\leanok
    For (1), $0\le\langle(K-H)v,v\rangle\le\langle Tv,v\rangle=0$, and a positive
    operator whose quadratic form vanishes at $v$ annihilates $v$. For (2),
    $\lVert(1-Q)v\rVert^2=\langle(1-Q)v,v\rangle\le2\langle Hv,v\rangle=0$. For (3),
    $0\le H\le K$ gives $\ker K\subseteq\ker H$ by
    Lemma~\ref{lem:gap_positive_comparison}. Conversely, if $Hv=0$, then $Qv=v$ by
    (2), and $Kv=KQv=HQv=Hv=0$.
\end{proof}

\begin{theorem}[A gap through a reducing projection]
    \label{thm:gap_reducing_projection}
    \lean{LinearMap.re_inner_ge_half_of_one_sub_le_twice,
        LinearMap.IsPositive.norm_gap_of_reducing_projection}
    \leanok
    Let $H,Q$ be operators on a complex inner product space with $1-Q\le2H$. Then
    $\tfrac12\lVert v\rVert^2\le\operatorname{Re}\langle Hv,v\rangle$ for every $v$ with
    $Qv=0$. Suppose in addition that the space is finite-dimensional, $H,K\ge0$, $Q$ is an
    orthogonal projection commuting with $H$, $H\le K$, $HQ=KQ$, and $\delta\ge0$ is a
    norm gap of $K$. Then $\min(\delta,1/2)$ is a norm gap of $H$.
\end{theorem}
\begin{proof}\leanok
    \uses{lem:gap_reducing_projection_kernel, lem:gap_energy_to_norm}
    If $Qv=0$, then $\lVert v\rVert^2=\langle(1-Q)v,v\rangle\le2\operatorname{Re}\langle
        Hv,v\rangle$. For the gap, $\ker H=\ker K$ by
    Lemma~\ref{lem:gap_reducing_projection_kernel}, and $Q$ fixes this kernel. Let
    $v\perp\ker H$ and write $v=x+y$ with $x=Qv$ and $y=(1-Q)v$. Then $x\perp\ker K$,
    since $\langle x,w\rangle=\langle v,Qw\rangle=\langle v,w\rangle=0$ for
    $w\in\ker K$, and the cross terms vanish because $Q$ commutes with $H$. Since
    $Hx=HQv=KQv=Kx$ and a norm gap $\delta\ge0$ of $K\ge0$ bounds its quadratic form
    on $(\ker K)^\perp$,
    \begin{align*}
        \operatorname{Re}\langle Hv,v\rangle
        =\operatorname{Re}\langle Kx,x\rangle+\operatorname{Re}\langle Hy,y\rangle
        \ge\delta\lVert x\rVert^2+\tfrac12\lVert y\rVert^2
        \ge\min(\delta,1/2)\lVert v\rVert^2,
    \end{align*}
    and Lemma~\ref{lem:gap_energy_to_norm} gives the norm gap.
\end{proof}

\subsection{Isometries}
\label{sec:martingale_gaps_isometries}

\begin{lemma}[Kernel and gap under an isometric conjugation]
    \label{lem:gap_isometric_conjugation}
    \lean{LinearIsometryEquiv.map_apply_eq_of_conj,
        LinearIsometryEquiv.mem_ker_iff_of_conj,
        LinearIsometryEquiv.ker_eq_map_of_conj,
        LinearIsometryEquiv.mem_orthogonal_ker_iff_of_conj,
        LinearIsometryEquiv.norm_gap_iff_of_conj}
    \leanok
    Let $U:E\to F$ be a surjective complex linear isometry of inner product spaces,
    and let $P$ on $E$ and $Q$ on $F$ satisfy $Q=UPU^{-1}$. Then $UPx=QUx$,
    $\ker Q=U(\ker P)$ and $(\ker Q)^\perp=U((\ker P)^\perp)$. For every $\delta\in\R$,
    $\delta$ is a norm gap of $P$ if and only if it is a norm gap of $Q$. Neither
    finite dimension nor positivity is needed.
\end{lemma}
\begin{proof}\leanok
    Substituting $Ux$ in the conjugation gives $UPx=QUx$. Injectivity of $U$
    identifies kernel membership, and surjectivity gives the kernel equality. The
    identity $\langle Uz,Ux\rangle=\langle z,x\rangle$ identifies the orthogonal
    complements, and applying the norm estimate at $x$ and at $Ux$, with
    $\lVert Ux\rVert=\lVert x\rVert$, gives the equivalence.
\end{proof}

\begin{lemma}[Projection onto an isometric image]
    \label{lem:gap_isometric_image_projection}
    \lean{LinearIsometry.starProjection_map_eq_comp_adjoint}
    \leanok
    Let $T:E\to F$ be an isometry between finite-dimensional real or complex inner
    product spaces and $G\subseteq E$ a subspace. Then $P_{TG}=TP_GT^*$.
\end{lemma}
\begin{proof}\leanok
    For $y\in F$, the vector $TP_GT^*y$ lies in $TG$, and for every $x\in G$ the
    adjoint identity and $T^*T=1$ give
    $\langle y-TP_GT^*y,Tx\rangle=\langle T^*y-P_GT^*y,x\rangle=0$. These two
    properties characterize the orthogonal projection onto $TG$.
\end{proof}

\begin{definition}[Isometric compression]
    \label{def:gap_compression}
    \lean{LinearIsometry.compression}
    \leanok
    For an isometry $U:E\to F$ of complex inner product spaces and an operator $H$ on
    $F$, the compression of $H$ is the operator $K=U^*HU$ on $E$.
\end{definition}

\begin{lemma}[Compression to a reducing image]
    \label{lem:gap_compression_intertwines}
    \lean{LinearIsometry.apply_compression_of_commute,
        LinearIsometry.compression_adjoint_apply_of_commute,
        LinearIsometry.ker_compression_of_commute,
        LinearIsometry.ker_compression_eq_map_adjoint_of_commute,
        LinearIsometry.map_mem_orthogonal_ker_compression_of_commute}
    \leanok
    \uses{def:gap_compression}
    Suppose that $UU^*H=HUU^*$, and let $K=U^*HU$. Then
    \begin{align*}
        UK                & =HU,                         & KU^* & =U^*H, \\
        \ker K            & =U^{-1}(\ker H)=U^*(\ker H), &
        U((\ker K)^\perp) & \subseteq(\ker H)^\perp.
    \end{align*}
    Neither surjectivity of $U$ nor symmetry of $H$ is needed.
\end{lemma}
\begin{proof}\leanok
    Combining $U^*U=1$ with the commutation gives $UK=UU^*HU=HU$ and
    $KU^*=U^*HUU^*=U^*H$. Injectivity of $U$ gives $\ker K=U^{-1}(\ker H)$. The second
    identity maps $\ker H$ into $\ker K$, and $v=U^*(Uv)$ gives the reverse inclusion
    in $U^*(\ker H)$. If $v\perp\ker K$ and $w\in\ker H$, then
    $\langle w,Uv\rangle=\langle U^*w,v\rangle=0$.
\end{proof}

\begin{theorem}[A norm gap survives compression to a reducing image]
    \label{thm:gap_compression}
    \lean{LinearIsometry.norm_gap_compression_of_commute}
    \leanok
    \uses{def:gap_compression}
    Suppose that $UU^*H=HUU^*$ and that $\delta\in\R$ is a norm gap of $H$. Then
    $\delta$ is a norm gap of $U^*HU$. Either space may be zero. This is the
    restriction step of~\cite[Section~5, lines 1695--1777]{GarreRubio2023MPOSym}.
\end{theorem}
\begin{proof}\leanok
    \uses{lem:gap_compression_intertwines}
    Let $K=U^*HU$ and $v\perp\ker K$. By
    Lemma~\ref{lem:gap_compression_intertwines}, $Uv\perp\ker H$ and $HUv=UKv$, so
    $\delta\lVert v\rVert=\delta\lVert Uv\rVert\le\lVert HUv\rVert=\lVert Kv\rVert$.
\end{proof}

\begin{lemma}[Nested subspaces]
    \label{lem:gap_supported_projection}
    \lean{Submodule.starProjection_comp_starProjection_eq_right_of_le,
        Submodule.starProjection_comp_starProjection_eq_left_of_le,
        Submodule.commute_starProjection_one_sub_starProjection_of_le,
        Submodule.one_sub_starProjection_apply_of_mem_orthogonal_of_le,
        Submodule.one_sub_starProjection_le_one_sub_starProjection_of_le,
        Submodule.isPositive_one_sub_starProjection_sub_one_sub_starProjection_of_le,
        Submodule.re_inner_one_sub_starProjection_le_of_le,
        Submodule.norm_sq_starProjection_orthogonal_le_re_inner_one_sub_of_le}
    \leanok
    Let $S\subseteq F$ be subspaces of a finite-dimensional complex Hilbert space, and
    put $h=1-P_S$ and $Q=P_F$. Then $QP_S=P_SQ=P_S$, $Qh=hQ$ and $1-Q\le h$; also
    $hv=v$ for $v\in F^\perp$, and
    $\lVert P_{F^\perp}v\rVert^2\le\operatorname{Re}\langle hv,v\rangle$ for every $v$.
    Equivalently, $h-(1-Q)\ge0$, so the quadratic form of $1-Q$ is bounded by that
    of $h$ at every vector.
\end{lemma}
\begin{proof}\leanok
    Projection onto $F$ fixes $S$, which gives $QP_S=P_S$; taking adjoints gives
    $P_SQ=P_S$. Thus $Q$ commutes with $h$. The order $P_S\le Q$ gives $1-Q\le h$. If
    $v\in F^\perp$, then $P_Sv=0$ and $hv=v$. Finally
    $\operatorname{Re}\langle(1-Q)v,v\rangle=\lVert P_{F^\perp}v\rVert^2$.
\end{proof}

\begin{theorem}[Compression of a complementary projection]
    \label{thm:gap_supported_compression}
    \lean{LinearIsometry.starProjection_range_eq_comp_adjoint,
        LinearIsometry.commute_rangeProjection_one_sub_starProjection_of_le_range,
        LinearIsometry.map_adjoint_map_eq_of_le_range,
        LinearIsometry.compression_starProjection_of_le_range,
        LinearIsometry.compression_one_sub_starProjection_of_le_range,
        LinearIsometry.compression_one_sub_starProjection_eq_orthogonal_of_le_range,
        LinearIsometry.ker_compression_one_sub_starProjection_of_le_range,
        LinearIsometry.norm_gap_compression_one_sub_starProjection_of_le_range}
    \leanok
    \uses{def:gap_compression}
    Let $U:E\to F$ be an isometry between finite-dimensional complex Hilbert spaces,
    and let $S\subseteq\operatorname{ran}U$. Then $P_{\operatorname{ran}U}=UU^*$,
    $U(U^*S)=S$, $UU^*$ commutes with $1-P_S$, and
    \begin{align*}
        U^*P_SU                     & =P_{U^*S},                    \\
        U^*(1-P_S)U                 & =1-P_{U^*S}=P_{(U^*S)^\perp}, \\
        \ker\bigl(U^*(1-P_S)U\bigr) & =U^*S.
    \end{align*}
    Every norm gap of $1-P_S$ is a norm gap of $U^*(1-P_S)U$.
\end{theorem}
\begin{proof}\leanok
    \uses{lem:gap_supported_projection, lem:gap_isometric_image_projection,
        thm:gap_compression}
    Lemma~\ref{lem:gap_isometric_image_projection} with $G=E$ gives
    $P_{\operatorname{ran}U}=UU^*$, and Lemma~\ref{lem:gap_supported_projection} with
    $F=\operatorname{ran}U$ gives the commutation. The identities $U^*U=1$ and
    $S\subseteq\operatorname{ran}U$ give $U(U^*S)=S$. By
    Lemma~\ref{lem:gap_isometric_image_projection}, $P_S=P_{U(U^*S)}=UP_{U^*S}U^*$;
    multiplying by $U^*$ on the left and $U$ on the right gives $U^*P_SU=P_{U^*S}$,
    and its complement and kernel follow. Theorem~\ref{thm:gap_compression} carries
    the gap through the compression.
\end{proof}

Let $U:E\to F$ be an isometry of complex inner product spaces, $K$ an operator on $E$
and $H$ an operator on $F$ with $HU=UK$. Every $v\in F$ decomposes as
\begin{align}
    v & =Ux+y, & x & =U^*v, & U^*y & =0, & \lVert v\rVert^2 & =\lVert x\rVert^2+\lVert y\rVert^2.
    \label{eq:gap_inclusion_split}
\end{align}
Write $q_H(y)=\operatorname{Re}\langle Hy,y\rangle$, and likewise $q_K$.

\begin{lemma}[The kernel under an isometric inclusion]
    \label{lem:gap_inclusion_kernel}
    \lean{LinearIsometry.ker_eq_map_of_intertwines_of_inactive_bound}
    \leanok
    If $HU=UK$, $c>0$, and $c\lVert y\rVert^2\le q_H(y)$ whenever $U^*y=0$, then
    $\ker H=U(\ker K)$. No symmetry or positivity is needed.
\end{lemma}
\begin{proof}\leanok
    If $Hv=0$, the decomposition~(\ref{eq:gap_inclusion_split}) gives $Hy=-UKx$, so
    $q_H(y)=-\operatorname{Re}\langle Kx,U^*y\rangle=0$, and the strict bound forces
    $y=0$. Thus $v=Ux$, and $UKx=Hv=0$ gives $Kx=0$ by injectivity of $U$. Conversely,
    $Kx=0$ gives $HUx=UKx=0$.
\end{proof}

\begin{lemma}[Combining quadratic-form bounds]
    \label{lem:gap_inclusion_quadratic}
    \lean{LinearIsometry.re_inner_ge_of_intertwines}
    \leanok
    Suppose $H$ is symmetric, $HU=UK$, $\delta,c\in\R$, $\delta\lVert x\rVert^2\le
        q_K(x)$ for $x\perp\ker K$, and $c\lVert y\rVert^2\le q_H(y)$ whenever $U^*y=0$.
    Then $\min\{\delta,c\}\lVert v\rVert^2\le q_H(v)$ for $v\perp\ker H$. Neither
    constant needs to be nonnegative.
\end{lemma}
\begin{proof}\leanok
    For $v\perp\ker H$, the vector $x=U^*v$ lies in $(\ker K)^\perp$, because
    $U(\ker K)\subseteq\ker H$. The intertwining gives $\langle HUx,y\rangle=
        \langle UKx,y\rangle=0$, and symmetry gives $\langle Hy,Ux\rangle=0$. Therefore,
    by~(\ref{eq:gap_inclusion_split}),
    \begin{align*}
        q_H(v)=q_K(x)+q_H(y)\ge\delta\lVert x\rVert^2+c\lVert y\rVert^2
        \ge\min\{\delta,c\}\bigl(\lVert x\rVert^2+\lVert y\rVert^2\bigr)
        =\min\{\delta,c\}\lVert v\rVert^2.
    \end{align*}
\end{proof}

\begin{theorem}[A norm gap through an isometric inclusion]
    \label{thm:gap_inclusion_norm}
    \lean{LinearIsometry.norm_gap_of_intertwines}
    \leanok
    Suppose $H,K\ge0$, $HU=UK$, $\delta\ge0$ is a norm gap of $K$, and
    $c\lVert y\rVert^2\le q_H(y)$ whenever $U^*y=0$. Then $\min\{\delta,c\}$ is a norm
    gap of $H$. No sign of $c$ is assumed.
\end{theorem}
\begin{proof}\leanok
    \uses{lem:gap_inclusion_quadratic, lem:gap_energy_to_norm}
    For the positive operator $K$, the norm gap $\delta\ge0$ gives
    $q_K(x)\ge\delta\lVert x\rVert^2$ on $(\ker K)^\perp$.
    Lemma~\ref{lem:gap_inclusion_quadratic} and Lemma~\ref{lem:gap_energy_to_norm}
    give the claim.
\end{proof}

\subsection{Bounded congruences}
\label{sec:martingale_gaps_congruence}

\begin{lemma}[The kernel of a congruence]
    \label{lem:gap_congruence_kernel}
    \lean{LinearEquiv.ker_adjoint_congruence}
    \leanok
    Let $e:E\to F$ be an invertible linear map of finite-dimensional complex Hilbert
    spaces and $P$ an operator on $F$. Then $\ker(e^*Pe)=e^{-1}(\ker P)$. No
    positivity is needed.
\end{lemma}
\begin{proof}\leanok
    If $e^*Pev=0$, then $\langle Pev,Pev\rangle=\langle e^*Pev,e^{-1}Pev\rangle=0$,
    since $e^*$ is the adjoint of $e$ and $e(e^{-1}Pev)=Pev$; hence $Pev=0$. The
    converse is immediate.
\end{proof}

\begin{theorem}[A gap under a bounded congruence]
    \label{thm:gap_bounded_congruence}
    \lean{LinearMap.IsPositive.norm_gap_adjoint_congruence_of_inverse_bound,
        LinearMap.IsPositive.bounded_congruence_gap}
    \leanok
    \uses{lem:gap_congruence_kernel}
    Let $P\ge0$ on $F$ have norm gap $\delta\ge0$, and let $e:E\to F$ be invertible with
    $\lVert e^{-1}w\rVert\le C\lVert w\rVert$ for every $w$, where $C>0$. Then
    $\delta/C^2$ is a norm gap of $Q=e^*Pe$. If $\delta>0$, then in addition
    $\delta/C^2>0$, $Q\ge0$ and $\ker Q=e^{-1}(\ker P)$.
\end{theorem}
\begin{proof}\leanok
    \uses{lem:gap_congruence_kernel}
    The quadratic form of $Q$ is $\langle Qv,v\rangle=\langle Pev,ev\rangle$, so
    $Q\ge0$, and Lemma~\ref{lem:gap_congruence_kernel} gives the kernel. Let
    $v\perp\ker Q$, and split $ev=P_{\ker P}ev+w$ with $w\perp\ker P$. The vector
    $u=e^{-1}P_{\ker P}ev$ lies in $\ker Q$, so $v\perp u$ and $e^{-1}w=v-u$; hence
    $\lVert v\rVert\le\lVert v-u\rVert=\lVert e^{-1}w\rVert\le C\lVert w\rVert$.
    Since $Pev=Pw$,
    \begin{align*}
        \tfrac{\delta}{C^2}\lVert v\rVert^2\le\delta\lVert w\rVert^2
        \le\operatorname{Re}\langle Pw,w\rangle=\operatorname{Re}\langle Qv,v\rangle
        \le\lVert Qv\rVert\lVert v\rVert,
    \end{align*}
    where the second inequality is the quadratic form of the positive operator $P$ on
    $(\ker P)^\perp$. Dividing by $\lVert v\rVert$ gives the norm gap.
\end{proof}

\subsection{Uniform gaps on compact sets}
\label{sec:martingale_gaps_compact}

\begin{lemma}[Uniform norm lower bounds on a common tail]
    \label{lem:gap_compact_tail}
    \lean{IsCompact.exists_uniform_norm_lower_bound_on_tail}
    \leanok
    \uses{lem:compact_uniform_pos_nat_bounds}
    Let $S$ be a compact subset of a topological space. For each $N\in\N$, let
    $E_N,F_N$ be seminormed groups, and let $H_N(x):E_N\to F_N$ be maps and
    $V_N(x)\subseteq E_N$ sets of vectors, for every parameter $x$. Suppose every
    $x\in S$ has a neighborhood $U_x$, a constant $\delta_x>0$ and a threshold $N_x$ with
    $\delta_x\lVert v\rVert\le\lVert H_N(y)v\rVert$ whenever $y\in U_x$, $N\ge N_x$ and
    $v\in V_N(y)$. Then one $\delta>0$ and one $N_0$ give
    $\delta\lVert v\rVert\le\lVert H_N(x)v\rVert$ for every $x\in S$, $N\ge N_0$ and
    $v\in V_N(x)$. The maps need not be linear. This is the norm form of the
    compactness argument of~\cite[Appendix~A]{Schuch2011Phases}.
\end{lemma}
\begin{proof}\leanok
    \uses{lem:compact_uniform_pos_nat_bounds}
    Decreasing the constant preserves the inequality because $\lVert v\rVert\ge0$, and
    increasing the threshold restricts the indices; apply
    Lemma~\ref{lem:compact_uniform_pos_nat_bounds}.
\end{proof}

\begin{theorem}[A uniform norm gap on a compact parameter set]
    \label{thm:gap_compact_kernel}
    \lean{ContinuousLinearMap.exists_uniform_norm_gap_of_compact}
    \leanok
    Let $E$ be a finite-dimensional complex Hilbert space, and let $x\mapsto H_x$ be a
    continuous family of operators on $E$ such that the orthogonal projection onto
    $\ker H_x$ is continuous in $x$. For every compact parameter set $S$ there is
    $\delta>0$ that is a norm gap of $H_x$ for every $x\in S$. No positivity or
    symmetry is required. This is the step of~\cite[Appendix~A, lines
        2575--2578]{Schuch2011Phases} that makes a gap uniform over a compact set.
\end{theorem}
\begin{proof}\leanok
    \uses{lem:gap_orthogonal_kernel, lem:kernel_gap_perturbation,
        lem:compact_uniform_pos_nat_bounds}
    By Lemma~\ref{lem:gap_orthogonal_kernel}, each $H_x$ has a positive norm gap.
    By Lemma~\ref{lem:kernel_gap_perturbation}, half of it is a norm gap of $H_y$ for
    all $y$ near $x$, by continuity of the operator and of its kernel projection.
    Lemma~\ref{lem:compact_uniform_pos_nat_bounds} gives one positive constant on
    $S$.
\end{proof}

\subsection{Products of commuting projections}
\label{sec:martingale_gaps_products}

\begin{lemma}[The product of complementary commuting projections]
    \label{lem:gap_commuting_projection_product}
    \lean{LinearMap.range_listProd_one_sub_eq_ker_sum,
        LinearMap.isProj_ker_sum_listProd_one_sub}
    \leanok
    Let $P_0,\dots,P_{N-1}$ be pairwise commuting orthogonal projections on a real or
    complex inner product space. The ordered product
    $D=(1-P_0)(1-P_1)\cdots(1-P_{N-1})$ is idempotent, and
    $\operatorname{ran}D=\ker\sum_iP_i$.
\end{lemma}
\begin{proof}\leanok
    Each $P_i$ commutes with every factor of $D$, and $P_i(1-P_i)=0$, so $P_iD=0$;
    hence $\operatorname{ran}D\subseteq\ker\sum_iP_i$. Conversely, if $\sum_iP_iv=0$,
    then $\sum_i\lVert P_iv\rVert^2=0$, so $P_iv=0$ for every $i$, and $Dv=v$. A
    product of commuting idempotents is idempotent.
\end{proof}

\section{Nachtergaele's martingale method}
\label{sec:martingale_nachtergaele}

Nachtergaele's criterion~\cite[Theorem~2.1(i)]{Nachtergaele1996Spectral} bounds the
gap of a frustration-free Hamiltonian on an increasing sequence of volumes
$\Lambda_0=\emptyset\subset\Lambda_1\subset\cdots\subset\Lambda_N$ from three conditions on
local Hamiltonians $H_n$ supported on windows of $l+1$ consecutive increments. With
$G_n$ the projection onto the ground space of the volume $\Lambda_n$ and $Q_n$ the
projection onto the ground space of the window $H_n$, condition C1 compares the sum of
the window Hamiltonians with the full Hamiltonian, condition C2 is a local gap for each
window, and condition C3 bounds the overlap $\lVert Q_nE_n\rVert$ of the local and the
global ground spaces~\cite[lines 1030--1094]{Nachtergaele1996Spectral}. This section
treats the argument for an abstract decreasing family of projections.

\begin{definition}[Nested ground projections and martingale differences]
    \label{def:martingale_nested_projections}
    \lean{FrustrationFree.NestedGroundProjections,
        FrustrationFree.NestedGroundProjections.martingaleDifference}
    \leanok
    Let $\mathcal H$ be a complex inner product space. A family of nested ground
    projections is a sequence $(G_n)_{n\ge0}$ of orthogonal projections on
    $\mathcal H$ with $\operatorname{ran}G_n\subseteq\operatorname{ran}G_m$ whenever
    $m\le n$. Its $n$th martingale difference is $E_n=G_n-G_{n+1}$.
\end{definition}

\begin{lemma}[The martingale differences are orthogonal]
    \label{lem:martingale_differences}
    \lean{FrustrationFree.NestedGroundProjections.martingaleDifference_isSymmetricProjection,
        FrustrationFree.NestedGroundProjections.martingaleDifference_comp_eq_zero,
        FrustrationFree.NestedGroundProjections.martingaleDifference_comp,
        FrustrationFree.NestedGroundProjections.inner_martingaleDifference_eq_zero,
        FrustrationFree.NestedGroundProjections.sum_martingaleDifference,
        FrustrationFree.NestedGroundProjections.sum_martingaleDifference_eq_id}
    \leanok
    \uses{def:martingale_nested_projections}
    As asserted after the definition of $E_n$ in~\cite{Nachtergaele1996Spectral}, each
    $E_n$ is an orthogonal projection, $E_mE_n=\delta_{mn}E_m$, and
    $\langle E_mx,E_ny\rangle=0$ for $m\neq n$. For every $N$,
    $\sum_{n=0}^{N-1}E_n=G_0-G_N$. Under the conventions $G_0=1$ and $G_{N+1}=0$ of the
    source, $\sum_{n=0}^{N}E_n=1$.
\end{lemma}
\begin{proof}\leanok
    If $m\le n$, nestedness and orthogonality give $G_mG_n=G_nG_m=G_n$. Hence
    $E_n^2=(G_n-G_{n+1})^2=G_n-G_{n+1}=E_n$ and $E_mE_n=0$ for $m\neq n$, and
    $E_n^*=E_n$ because $G_n$ and $G_{n+1}$ are symmetric. The orthogonality of the
    images follows from $\langle E_mx,E_ny\rangle=\langle x,E_mE_ny\rangle$, and the
    sum telescopes.
\end{proof}

\begin{lemma}[Truncated reconstruction]
    \label{lem:martingale_reconstruction}
    \lean{FrustrationFree.NestedGroundProjections.sum_martingaleDifference_apply_of_mem_orthogonal}
    \leanok
    \uses{def:martingale_nested_projections}
    Suppose $G_0=1$. As in the identity \textup{resolutionpsi} in the proof
    of~\cite[Theorem~2.1(i)]{Nachtergaele1996Spectral}, every
    $v\perp\operatorname{ran}G_N$ satisfies $v=\sum_{n=0}^{N-1}E_nv$.
\end{lemma}
\begin{proof}\leanok
    \uses{lem:martingale_differences}
    Since $v\perp\operatorname{ran}G_N$, $G_Nv=0$, and
    Lemma~\ref{lem:martingale_differences} gives $\sum_{n=0}^{N-1}E_nv=(G_0-G_N)v=v$.
\end{proof}

\begin{lemma}[The Pythagorean identity for martingale differences]
    \label{lem:martingale_pythagorean}
    \lean{FrustrationFree.NestedGroundProjections.norm_sq_sum_martingaleDifference_finset,
        FrustrationFree.NestedGroundProjections.norm_sq_sum_martingaleDifference}
    \leanok
    \uses{def:martingale_nested_projections}
    For every $v\in\mathcal H$ and finite $S\subset\N$,
    $\bigl\lVert\sum_{n\in S}E_nv\bigr\rVert^2=\sum_{n\in S}\lVert E_nv\rVert^2$; in
    particular this holds for $S=\{0,\dots,N-1\}$.
\end{lemma}
\begin{proof}\leanok
    \uses{lem:martingale_differences}
    The summands are pairwise orthogonal by Lemma~\ref{lem:martingale_differences},
    and the finite Pythagorean theorem applies.
\end{proof}

\begin{lemma}[The norm on the final orthogonal complement]
    \label{lem:martingale_norm_decomposition}
    \lean{FrustrationFree.NestedGroundProjections.norm_sq_eq_sum_martingaleDifference_of_mem_orthogonal}
    \leanok
    \uses{def:martingale_nested_projections}
    Suppose $G_0=1$. As in the norm identity that follows \textup{resolutionpsi} in the
    proof of~\cite[Theorem~2.1(i)]{Nachtergaele1996Spectral}, every
    $v\perp\operatorname{ran}G_N$ satisfies
    $\lVert v\rVert^2=\sum_{n=0}^{N-1}\lVert E_nv\rVert^2$.
\end{lemma}
\begin{proof}\leanok
    \uses{lem:martingale_reconstruction, lem:martingale_pythagorean}
    Substitute $v=\sum_{n=0}^{N-1}E_nv$ from Lemma~\ref{lem:martingale_reconstruction}
    into Lemma~\ref{lem:martingale_pythagorean}.
\end{proof}

\begin{lemma}[Martingale differences below a fixed level]
    \label{lem:martingale_below_level}
    \lean{FrustrationFree.NestedGroundProjections.martingaleDifference_apply_eq_zero_of_lt}
    \leanok
    \uses{def:martingale_nested_projections}
    If $G_{n_0}v=v$, then $E_nv=0$ for every $n<n_0$.
\end{lemma}
\begin{proof}\leanok
    For $m\le n_0$ the ranges are nested, so $G_mG_{n_0}=G_{n_0}$ and
    $G_mv=G_mG_{n_0}v=v$. For $n<n_0$ this applies to $n$ and to $n+1$, so
    $E_nv=v-v=0$.
\end{proof}

\begin{lemma}[Cancellation outside the window]
    \label{lem:martingale_outside_window}
    \lean{FrustrationFree.NestedGroundProjections.martingaleDifference_comp_localProjection_comp_eq_zero,
        FrustrationFree.NestedGroundProjections.sum_martingaleDifference_localProjection_eq_sum_Icc,
        FrustrationFree.NestedGroundProjections.inner_sum_martingaleDifference_localProjection_eq_inner_sum_Icc}
    \leanok
    \uses{def:martingale_nested_projections}
    Let $Q$ be an operator with $E_mQ=QE_m$ whenever $m<n-l$ or $n<m$. Then
    $E_mQE_n=0$ for every such $m$. If $n<N$, then
    \begin{align}
        \sum_{m=0}^{N-1}E_mQE_nv & =\sum_{m=n-l}^{n}E_mQE_nv,
        \notag                                                                          \\
        \Bigl\langle v,\sum_{m=0}^{N-1}E_mQE_nv\Bigr\rangle
                                 & =\Bigl\langle\sum_{m=n-l}^{n}E_mv,QE_nv\Bigr\rangle,
        \label{eq:martingale_outside_window}
    \end{align}
    where $n-l$ is read as $0$ when $l>n$. This is the reduction from \textup{Enpsi} to
    the first line of \textup{Enpsi2} in~\cite[lines
        1206--1220]{Nachtergaele1996Spectral}.
\end{lemma}
\begin{proof}\leanok
    \uses{lem:martingale_differences}
    For an index outside the window, $E_mQE_n=QE_mE_n=0$ by
    Lemma~\ref{lem:martingale_differences}, since $m\neq n$. Hence only
    $n-l\le m\le n$ remains, and $n<N$ places these indices in the range of the sum.
    Symmetry of $E_m$ gives $\langle v,E_mQE_nv\rangle=\langle E_mv,QE_nv\rangle$.
\end{proof}

\begin{lemma}[Weighted inner-product estimate]
    \label{lem:martingale_weighted_inner}
    \lean{FrustrationFree.re_inner_le_weighted_norm_sq}
    \leanok
    For $x,y$ in a complex inner product space and $c>0$,
    $\operatorname{Re}\langle x,y\rangle\le\frac1{2c}\lVert x\rVert^2+\frac c2\lVert y\rVert^2$.
    This is the estimate used twice in the display \textup{Enpsi2}
    of~\cite[lines 1223--1239]{Nachtergaele1996Spectral}.
\end{lemma}
\begin{proof}\leanok
    The Cauchy--Schwarz inequality bounds the left side by
    $\lVert x\rVert\lVert y\rVert$, and the weighted arithmetic--geometric mean
    inequality gives the bound.
\end{proof}

\begin{lemma}[The overlap estimate on the range of a projection]
    \label{lem:martingale_c3_projection}
    \lean{FrustrationFree.norm_sq_apply_projection_le_of_norm_comp_le}
    \leanok
    Let $P$ be an orthogonal projection on a complex inner product space and $Q$ a
    bounded operator from that space to another with $\lVert QP\rVert\le\epsilon$.
    Then $\lVert QPv\rVert^2\le\epsilon^2\lVert Pv\rVert^2$ for every $v$; the case
    $\epsilon=0$ is included. This is the form of condition C3
    of~\cite[lines 1083--1094]{Nachtergaele1996Spectral} used at lines 1240--1249.
\end{lemma}
\begin{proof}\leanok
    Idempotency gives $QPv=QP(Pv)$. The operator-norm bound at $Pv$, squared, gives
    the claim; the hypothesis forces $\epsilon\ge0$.
\end{proof}

\begin{lemma}[Counting moving windows]
    \label{lem:martingale_moving_window}
    \lean{FrustrationFree.movingWindow_sum_le, FrustrationFree.movingWindow_sum_range_le,
        FrustrationFree.movingWindow_sum_Ico_le}
    \leanok
    Let $(a_m)_{m\ge0}$ be nonnegative reals, $l,n_0,N\in\N$, and read $n-l$ as $0$ when
    $l>n$. Then
    \begin{align}
        \sum_{n=0}^{N-1}\sum_{m=n-l}^{n}a_m   & \le(l+1)\sum_{m=0}^{N-1}a_m,
        \label{eq:martingale_moving_window}                                      \\
        \sum_{n=n_0}^{N-1}\sum_{m=n-l}^{n}a_m & \le(l+1)\sum_{m=n_0-l}^{N-1}a_m,
        \label{eq:martingale_moving_window_threshold}
    \end{align}
    and~(\ref{eq:martingale_moving_window}) holds as well when the terms with $n<l$
    are omitted on the left.
\end{lemma}
\begin{proof}\leanok
    With $k=n-m$, each pair in the double sum has $0\le k\le l$. For fixed $k$, the map
    $n\mapsto n-k$ is injective with image in $\{0,\dots,N-1\}$, so nonnegativity bounds
    its contribution by $\sum_{m=0}^{N-1}a_m$; summing over the $l+1$ values of $k$
    gives~(\ref{eq:martingale_moving_window}). This is the count after \textup{Enpsi2}
    in the proof of~\cite[Theorem~2.1(i), lines 1220--1255]{Nachtergaele1996Spectral}:
    for $a_m=\lVert E_m\psi\rVert^2$, each term occurs in at most $l+1$ windows. For
    (\ref{eq:martingale_moving_window_threshold}), replace $a_m$ by the sequence that
    vanishes below $n_0-l$ and agrees with $a_m$ from there on; every window with
    $n\ge n_0$ is unchanged, and~(\ref{eq:martingale_moving_window}) applies.
\end{proof}

\begin{theorem}[The energy estimate above a threshold]
    \label{thm:martingale_threshold_estimate}
    \lean{FrustrationFree.NestedGroundProjections.energy_lower_bound_of_nachtergaele_c1_c3_threshold}
    \leanok
    \uses{def:martingale_nested_projections}
    Let $(G_n)$ be nested ground projections on a finite-dimensional complex Hilbert
    space $\mathcal H$ with $G_0=1$, and fix $l,N,n_0\in\N$ with $n_0\le N$. Let $Q_n$ be
    orthogonal projections and $H_n$ operators for $n_0\le n<N$, and let $H$ be an
    operator, such that $E_mQ_n=Q_nE_m$ whenever $m<n-l$ or $n<m$. Write
    $q_n(x)=\operatorname{Re}\langle H_nx,x\rangle$ and
    $q(x)=\operatorname{Re}\langle Hx,x\rangle$. Assume, for every $x\in\mathcal H$ and
    $n_0\le n<N$, the quadratic-form versions of conditions C1 and C2 and condition C3,
    \begin{align}
        0                                  & \le\sum_{n=n_0}^{N-1}q_n(x)\le d_{l+1}q(x),
        \label{eq:martingale_threshold_c1}                                               \\
        \gamma_{l+1}\lVert(1-Q_n)x\rVert^2 & \le q_n(x),
        \label{eq:martingale_threshold_c2}                                               \\
        \lVert Q_nE_n\rVert                & \le\epsilon_l,
        \label{eq:martingale_threshold_c3}
    \end{align}
    with $\gamma_{l+1},d_{l+1}>0$ and $0\le\epsilon_l<1/\sqrt{l+1}$. Put
    $\kappa=(\gamma_{l+1}/d_{l+1})(1-\epsilon_l\sqrt{l+1})^2$. Then every
    $v\perp\operatorname{ran}G_N$ satisfies
    \begin{align}
        \kappa\sum_{n=n_0}^{N-1}\lVert E_nv\rVert^2
        \le\operatorname{Re}\langle Hv,v\rangle
        +\frac{\gamma_{l+1}}{d_{l+1}}\bigl(1-\epsilon_l\sqrt{l+1}\bigr)\epsilon_l\sqrt{l+1}
        \sum_{m=n_0-l}^{n_0-1}\lVert E_mv\rVert^2.
        \label{eq:martingale_threshold_estimate}
    \end{align}
    Conditions C2 and C3 have the form of~\cite[lines
        1043--1058 and 1083--1094]{Nachtergaele1996Spectral}, with their onset represented
    by $n_0$. The inequality~(\ref{eq:martingale_threshold_c1}) follows
    from the source condition C1 when $l\le n_0$ and the omitted local terms are
    nonnegative; for $n_0<l$ it is a stronger assumption. The coefficient $\kappa$ is
    the one of~\cite[Theorem~2.1(i)]{Nachtergaele1996Spectral}, but it multiplies the
    martingale mass above the threshold rather than $\lVert v\rVert^2$, and the
    correction bounds the contribution of the at most $l$ martingale differences just
    below the threshold. The source concludes the bound on $\lVert v\rVert^2$; the two
    agree exactly when the martingale differences below the threshold annihilate $v$,
    and the printed proof gives no argument for those indices. This scope restriction
    is recorded in~\cite{gap:nachtergaele96_theorem_2_1_lower_endpoint}.
\end{theorem}
\begin{proof}\leanok
    \uses{lem:martingale_reconstruction, lem:martingale_pythagorean,
        lem:martingale_outside_window, lem:martingale_weighted_inner,
        lem:martingale_c3_projection, lem:martingale_moving_window}
    Put $P=\sum_{n=n_0}^{N-1}\lVert E_nv\rVert^2$ and
    $S=\sum_{m=n_0-l}^{n_0-1}\lVert E_mv\rVert^2$. The reconstruction
    $v=\sum_{n=0}^{N-1}E_nv$ of Lemma~\ref{lem:martingale_reconstruction} and the
    cancellation~(\ref{eq:martingale_outside_window}) reduce the identity
    \textup{Enpsi} at each index $n_0\le n<N$ to the window $n-l\le m\le n$. Apply
    Lemma~\ref{lem:martingale_weighted_inner} with parameters $c_1,c_2>0$, then C2 and,
    through Lemma~\ref{lem:martingale_c3_projection}, C3, and sum over
    $n_0\le n<N$. Lemma~\ref{lem:martingale_pythagorean} splits each window norm into
    its summands, and~(\ref{eq:martingale_moving_window_threshold}) bounds the double
    sum by $(l+1)(S+P)$. With C1 this gives
    \begin{align*}
        \Bigl(2-c_1-\frac{\epsilon_l^2}{c_2}-(l+1)c_2\Bigr)P
        \le\frac{d_{l+1}}{c_1\gamma_{l+1}}\operatorname{Re}\langle Hv,v\rangle+(l+1)c_2S.
    \end{align*}
    For $\epsilon_l>0$, the choice $c_1=1-\epsilon_l\sqrt{l+1}$ and
    $c_2=\epsilon_l/\sqrt{l+1}$ gives $\epsilon_l^2/c_2=(l+1)c_2=\epsilon_l\sqrt{l+1}$
    and a left factor $1-\epsilon_l\sqrt{l+1}$; multiplying by
    $c_1\gamma_{l+1}/d_{l+1}$ gives~(\ref{eq:martingale_threshold_estimate}). For
    $\epsilon_l=0$ the choice $c_2=0$ is not admissible; instead C3 gives $Q_nE_n=0$, so
    the second inner product in \textup{Enpsi} vanishes, and the first weighted estimate
    with $c_1=1$ gives $P\le(d_{l+1}/\gamma_{l+1})\operatorname{Re}\langle Hv,v\rangle$,
    which is~(\ref{eq:martingale_threshold_estimate}) with vanishing correction. This
    local fix of the zero constant is recorded
    in~\cite{gap:nachtergaele96_theorem_2_1_lower_endpoint}.
\end{proof}

\begin{theorem}[Nachtergaele's estimate with no mass below the threshold]
    \label{thm:martingale_below_threshold_estimate}
    \lean{FrustrationFree.NestedGroundProjections.energy_lower_bound_of_nachtergaele_c1_c3_of_martingaleDifference_below_eq_zero}
    \leanok
    \uses{thm:martingale_threshold_estimate}
    Assume the setting of Theorem~\ref{thm:martingale_threshold_estimate} with
    $l\le n_0$, but replace~(\ref{eq:martingale_threshold_c1}) by condition C1 on the range
    $l\le n<M$ of its own window sum at the two volumes at which the source asserts it:
    $0\le\sum_{n=l}^{N-1}q_n(x)\le d_{l+1}q(x)$ at $M=N$, and
    $0\le\sum_{n=l}^{n_0-1}q_n(x)$ at $M=n_0$, for every $x$. Assume
    in addition that $E_nv=0$ for every $n<n_0$. Then
    $\kappa\lVert v\rVert^2\le\operatorname{Re}\langle Hv,v\rangle$, which is the
    conclusion of~\cite[Theorem~2.1(i), lines 1119--1130]{Nachtergaele1996Spectral} with
    each of C1, C2 and C3 on its source range.

    The added vanishing hypothesis is a local fix of the source theorem, recorded
    in~\cite{gap:nachtergaele96_theorem_2_1_lower_endpoint}, and it is the minimal one:
    the two below-threshold quantities of~(\ref{eq:martingale_threshold_estimate}) enter
    the conclusion with nonnegative weights, so both must vanish, and the vanishing of
    the omitted martingale mass forces the vanishing of the correction. Without it the
    statement is false, by Theorem~\ref{thm:martingale_lower_endpoint}.
\end{theorem}
\begin{proof}\leanok
    \uses{thm:martingale_threshold_estimate, lem:martingale_norm_decomposition}
    Condition C1 on $n_0\le n<N$ follows: its upper bound by subtracting the sum over
    $l\le n<n_0$, nonnegative by C1 at the volume $\Lambda_{n_0}$, and its nonnegativity
    termwise from C2. The correction in~(\ref{eq:martingale_threshold_estimate}) is a
    sum of squared norms of martingale differences with index below $n_0$, hence zero.
    Lemma~\ref{lem:martingale_norm_decomposition} splits $\lVert v\rVert^2$ into the sum
    over $n<n_0$, which is zero for the same reason, and the martingale mass above the
    threshold, which is therefore $\lVert v\rVert^2$.
\end{proof}

\begin{theorem}[The source conditions on their own ranges do not give the estimate]
    \label{thm:martingale_lower_endpoint}
    \lean{FrustrationFree.UnrestrictedNachtergaeleEstimate,
        FrustrationFree.not_unrestrictedNachtergaeleEstimate}
    \leanok
    \uses{thm:martingale_below_threshold_estimate}
    In the notation of Theorem~\ref{thm:martingale_below_threshold_estimate}: for every
    $k$ there are, on a complex Hilbert space of dimension $k+2$, nested ground
    projections $(G_n)$ with $G_0=1$, $G_{N+1}=0$ and $G_n\neq0$ for $n\le N$, nonzero
    orthogonal projections $Q_n$, operators $H_n$, volume Hamiltonians $H_{\Lambda_M}$,
    and a vector $v\perp\operatorname{ran}G_N$ with the following properties, for $l=0$,
    $n_0=1$, $N=k+1$, $\gamma_1=d_1=1$ and $\epsilon_0=0$. The commutations hold;
    condition C1 holds with the one constant $d_1$ at every volume $\Lambda_M$ on the
    range $l\le n<M$ of its own window sum; the gap inequality of C2 and condition C3
    hold on $n_0\le n<N$; every ground space is the kernel of its own Hamiltonian, that is
    $\ker H_{\Lambda_M}=\operatorname{ran}G_M$ for $M\le N$ and
    $\ker H_n=\operatorname{ran}Q_n$ for $0\le n<N$; and
    \begin{align}
        \operatorname{Re}\langle H_{\Lambda_N}v,v\rangle=\tfrac12\lVert v\rVert^2
        <\frac{\gamma_1}{d_1}(1-\epsilon_0)^2\lVert v\rVert^2.
        \label{eq:martingale_lower_endpoint}
    \end{align}
    Hence condition C1 at every volume on the range of its own window sum, together
    with C2 and C3 on $n_0\le n<N$, the ranges at which~\cite{Nachtergaele1996Spectral}
    imposes them, do not imply the conclusion of its Theorem~2.1(i), even when every
    ground space is the nonzero kernel of its Hamiltonian and under the conventions of
    the source at the extreme volumes. This is recorded
    in~\cite{gap:nachtergaele96_theorem_2_1_lower_endpoint}.
\end{theorem}
\begin{proof}\leanok
    Fix an orthonormal basis $e_0,\dots,e_{k+1}$ in which every operator is diagonal;
    the first $N=k+1$ vectors carry the single excitations of a chain of $N$ sites and
    the last its ground state. Let $G_n$ have weight $1$ at $e_i$ when $n\le i$ and $0$
    otherwise. Then $G_0=1$, $G_{N+1}=0$, every $G_n$ with $n\le N$ retains $e_{k+1}$,
    $G_N$ is the projection onto $e_{k+1}$, and $E_n$ is the projection onto $e_n$ for
    $n<N$. Let $c_0=\tfrac12$, $c_i=1$ for $1\le i\le k$ and $c_{k+1}=0$; let
    $H_{\Lambda_M}$ have weight $c_i$ at $e_i$ when $i<M$ and $0$ otherwise, let $H_n$
    have weight $c_n$ at $e_n$ and $0$ elsewhere, and let $Q_n=1-E_n$, nonzero since the
    dimension is at least two. Diagonality gives the commutations, and $Q_nE_n=0$ gives
    C3 with $\epsilon_0=0$. The sum of the $H_n$ over $0\le n<M$ is $H_{\Lambda_M}$, and
    the weights are nonnegative, so C1 holds at every volume with $d_1=1$. Since
    $c_i\neq0$ for $i\le k$, the kernel of $H_{\Lambda_M}$ is spanned by the $e_i$ with
    $i\ge M$ when $M\le N$, which is $\operatorname{ran}G_M$, and the kernel of $H_n$
    is spanned by the $e_i$ with $i\neq n$, which is $\operatorname{ran}Q_n$. For
    $1\le n<N$, $c_n=1$, so $\lVert(1-Q_n)x\rVert^2=\lVert E_nx\rVert^2=
        \operatorname{Re}\langle H_nx,x\rangle$, which is C2 with $\gamma_1=1$. The vector
    $v=e_0$ is orthogonal to $\operatorname{ran}G_N=\C e_{k+1}$, and
    $\operatorname{Re}\langle H_{\Lambda_N}v,v\rangle=c_0=\tfrac12$ while
    $\lVert v\rVert=1$, which is~(\ref{eq:martingale_lower_endpoint}). The whole
    martingale mass of $v$ sits at $n=0$, which C1 covers while C2 and C3 exclude it.
\end{proof}

\begin{theorem}[The energy estimate on the full range]
    \label{thm:martingale_full_range_estimate}
    \lean{FrustrationFree.NestedGroundProjections.energy_lower_bound_of_nachtergaele_c1_c3_full_range}
    \leanok
    \uses{def:martingale_nested_projections}
    Assume the setting of Theorem~\ref{thm:martingale_threshold_estimate} with $n_0=0$,
    that is, with~(\ref{eq:martingale_threshold_c1})--(\ref{eq:martingale_threshold_c3})
    for every $0\le n<N$. Then every $v\perp\operatorname{ran}G_N$ satisfies
    $\kappa\lVert v\rVert^2\le\operatorname{Re}\langle Hv,v\rangle$, with the
    coefficient $\kappa=(\gamma_{l+1}/d_{l+1})(1-\epsilon_l\sqrt{l+1})^2$
    of~\cite[Theorem~2.1(i)]{Nachtergaele1996Spectral}. This is a scope restriction of
    the source theorem, recorded
    in~\cite{gap:nachtergaele96_theorem_2_1_lower_endpoint}: it assumes C1, C2 and C3 on
    the full range $0\le n<N$, whereas the source imposes lower thresholds on them, and
    by Theorem~\ref{thm:martingale_lower_endpoint} those thresholds cannot simply be
    dropped.
\end{theorem}
\begin{proof}\leanok
    \uses{thm:martingale_threshold_estimate, lem:martingale_norm_decomposition}
    In Theorem~\ref{thm:martingale_threshold_estimate} with $n_0=0$, the correction is
    a sum over an empty range, and the martingale mass over $0\le n<N$ is
    $\lVert v\rVert^2$ by Lemma~\ref{lem:martingale_norm_decomposition}.
\end{proof}

\begin{theorem}[A norm gap from the full-range estimate]
    \label{thm:martingale_full_range_norm_gap}
    \lean{FrustrationFree.NestedGroundProjections.norm_lower_bound_of_nachtergaele_c1_c3_full_range}
    \leanok
    \uses{thm:martingale_full_range_estimate}
    In the setting of Theorem~\ref{thm:martingale_full_range_estimate}, suppose $H\ge0$
    and $\operatorname{ran}G_N=\ker H$. Then $\kappa$ is a norm gap of $H$.
\end{theorem}
\begin{proof}\leanok
    \uses{thm:martingale_full_range_estimate, lem:gap_energy_to_norm,
        lem:gap_quadratic_form_to_norm}
    For any $x$, let $w$ be $x$ minus its projection onto $\ker H$. Then $Hx=Hw$ and
    $\operatorname{Re}\langle Hx,x\rangle=\operatorname{Re}\langle Hw,w\rangle$.
    Theorem~\ref{thm:martingale_full_range_estimate} and
    Lemma~\ref{lem:gap_energy_to_norm} give $\kappa\lVert w\rVert\le\lVert Hx\rVert$, so
    $\kappa\operatorname{Re}\langle Hx,x\rangle\le\kappa\lVert Hx\rVert\lVert w\rVert
        \le\lVert Hx\rVert^2$, and Lemma~\ref{lem:gap_quadratic_form_to_norm} gives the
    norm gap.
\end{proof}

\begin{theorem}[A sharp gap from two ground spaces]
    \label{thm:martingale_two_ground_projections}
    \lean{FrustrationFree.norm_gap_of_two_ground_projection_bounds}
    \leanok
    Let $H$ be an operator on a finite-dimensional complex Hilbert space, $W=\ker H$,
    and $W\subseteq U\cap V$ for subspaces $U,V$. Suppose $\kappa\ge0$, $c<1$,
    $\lVert P_W-P_UP_V\rVert\le c$ and $\kappa\bigl((1-P_U)+(1-P_V)\bigr)\le H$. Then
    $\kappa(1-c)$ is a norm gap of $H$. This is a two-subspace consequence of the
    ground-projection estimate of~\cite[Section~6, the three-interval commutation lemma,
        part~(ii)]{Nachtergaele1996Spectral}.
\end{theorem}
\begin{proof}\leanok
    \uses{lem:gap_quadratic_form_to_norm}
    If $z\in U\cap V$ and $z\perp W$, then $\lVert z\rVert=\lVert(P_W-P_UP_V)z\rVert\le
        c\lVert z\rVert$, hence $z=0$; subtracting $P_Wv$ from $v\in U\cap V$ shows
    $W=U\cap V$. Put $L=(1-P_U)+(1-P_V)$. The projection defect bounds the cosine of the
    Friedrichs angle between $U$ and $V$ by $c$, and the anticommutator inequality for
    the complementary projections gives $L^2\ge(1-c)L$, with $\ker L=U\cap V=W$.
    Lemma~\ref{lem:gap_quadratic_form_to_norm} gives
    $(1-c)\lVert v\rVert^2\le\langle Lv,v\rangle$ for $v\perp W$, and $\kappa L\le H$
    with the Cauchy--Schwarz inequality gives
    $\kappa(1-c)\lVert v\rVert^2\le\langle Hv,v\rangle\le\lVert Hv\rVert\lVert v\rVert$.
\end{proof}

\section{Sweeps of projections}
\label{sec:martingale_sweep}

\begin{theorem}[One sweep of projections contracts]
    \label{thm:martingale_sweep_contraction}
    \lean{ContinuousLinearMap.complementProd_apply_of_forall_apply_eq_zero,
        ContinuousLinearMap.inner_complementProd_of_forall_apply_eq_zero,
        ContinuousLinearMap.norm_sq_complementProd_le_of_gap,
        ContinuousLinearMap.norm_pow_complementProd_le_of_gap}
    \leanok
    Let $\Pi_0,\dots,\Pi_{m-1}$ be orthogonal projections on a real or complex Hilbert
    space, each failing to commute with at most $d$ of the others, and let $\Omega$ be a
    vector with $\Pi_i\Omega=0$ for every $i$. Suppose
    $\gamma\lVert v\rVert^2\le\operatorname{Re}\langle v,\sum_i\Pi_iv\rangle$ for every
    $v\perp\Omega$. Put $D=(1-\Pi_{m-1})\cdots(1-\Pi_0)$. Then $D\Omega=\Omega$,
    $\langle\Omega,Dv\rangle=\langle\Omega,v\rangle$ for every $v$, and for $v\perp\Omega$
    \begin{align*}
        \bigl((d+1)^2+\gamma\bigr)\lVert Dv\rVert^2 & \le(d+1)^2\lVert v\rVert^2, \\
        \lVert D^tv\rVert                           & \le\theta^t\lVert v\rVert,
    \end{align*}
    where $\theta=(1+\gamma/(d+1)^2)^{-1/2}$ and the second bound assumes $\gamma\ge0$.
    This is the spectral gap amplification of~\cite{Anshu2016Simple}.
\end{theorem}
\begin{proof}\leanok
    \uses{thm:detectability_lemma}
    The detectability lemma, Theorem~\ref{thm:detectability_lemma}, gives
    $\langle Dv,\sum_i\Pi_iDv\rangle\le d^2(\lVert v\rVert^2-\lVert Dv\rVert^2)$. Each
    factor $1-\Pi_i$ is symmetric and fixes $\Omega$, so $D\Omega=\Omega$, $D^*\Omega=
        \Omega$, and $Dv\perp\Omega$; the left side is therefore at least
    $\gamma\lVert Dv\rVert^2$. Adding $(2d+1)\lVert Dv\rVert^2\le(2d+1)\lVert v\rVert^2$,
    which holds because $D$ is a product of orthogonal projections, gives the first
    bound. Since $D$ preserves the orthogonal complement of $\Omega$, the second bound
    follows by induction on $t$.
\end{proof}

\begin{lemma}[Range projections of positive local terms]
    \label{lem:martingale_range_projection}
    \lean{Matrix.PosSemidef.le_smul_supportProj, ManyBody.supportProj_mem_supportedOperators}
    \leanok
    \uses{def:sites_supported}
    For a positive semidefinite matrix $A$, let $\Pi_A$ be the orthogonal projection
    onto the range of $A$, so that $\Pi_Av=0$ whenever $Av=0$. If $A\le J$, then
    $A\le J\Pi_A$. If a Hermitian operator $A$ on a system of sites is supported on a
    region $D$, so is $\Pi_A$.
\end{lemma}
\begin{proof}\leanok
    \uses{thm:support_projection_order_kernel, cor:sites_cfc_support}
    By Theorem~\ref{thm:support_projection_order_kernel}, $\Pi_A=f(A)$ with $f(0)=0$ and
    $f(x)=1$ for $x\neq0$. On the spectrum of $A$, which lies in $[0,J]$,
    $x\le Jf(x)$, so $A\le J\Pi_A$. Corollary~\ref{cor:sites_cfc_support} shows that
    $f(A)$ is supported on $D$.
\end{proof}

\section{Commuting parent Hamiltonians on three regions}
\label{sec:martingale_commuting_parent}

Let $a$ and $b$ be disjoint regions of a finite system of sites, and let
$m=V\setminus(a\cup b)$ be the remaining sites, so that $V\setminus b=a\cup m$ and
$V\setminus a=m\cup b$. This section restates~\cite[Appendix~D.2, lines
    2187--2289]{Cirac2016MPDO_arXiv} for an interaction on the system.

\begin{definition}[Tripartite coordinates]
    \label{def:commuting_parent_coordinates}
    \lean{ManyBody.tripartiteEquiv, ManyBody.complRightEquiv, ManyBody.complLeftEquiv}
    \leanok
    \uses{def:sites_config}
    The tripartite coordinates are the bijection $X_V\to X_a\times X_m\times X_b$,
    $\sigma\mapsto(\sigma|_a,\sigma|_m,\sigma|_b)$, which identifies $\mathcal H_X$ with
    $\mathcal H_a\otimes\mathcal H_m\otimes\mathcal H_b$, together with the bijections
    $X_{V\setminus b}\to X_a\times X_m$ and $X_{V\setminus a}\to X_m\times X_b$ given by
    restriction.
\end{definition}

\begin{lemma}[Placements in product coordinates]
    \label{lem:commuting_parent_placement}
    \lean{ManyBody.reindex_place_eq_kronecker_one, ManyBody.reindex_place_eq_one_kronecker,
        ManyBody.reindex_place_left, ManyBody.reindex_place_right,
        ManyBody.reindex_place_compl_right, ManyBody.reindex_place_compl_left}
    \leanok
    \uses{def:sites_place, def:commuting_parent_coordinates}
    Let $g:X_V\to R\times C$ be a bijection whose first coordinate is $g_1(\sigma|_D)$ for
    a bijection $g_1:X_D\to R$, and whose second coordinates of $\sigma$ and $\tau$ agree
    exactly when $\sigma$ and $\tau$ agree outside $D$. Then the placement of a matrix
    $K$ on $D$ reads $K\otimes1$ in the coordinates $g$, with $K$ read through $g_1$; the
    same holds with the two factors exchanged. In the tripartite coordinates, the
    placements of $K$ on $a$, on $b$, on $V\setminus b$ and on $V\setminus a$ are
    $K\otimes1\otimes1$, $1\otimes1\otimes K$, $K\otimes1$ and $1\otimes K$, with $K$
    read in the coordinates $X_a\times X_m$ and $X_m\times X_b$ in the last two cases.
\end{lemma}
\begin{proof}\leanok
    Compare matrix elements with Definition~\ref{def:sites_place}: the element of the
    placement at $(\sigma,\tau)$ is $K_{\sigma|_D,\tau|_D}$ when $\sigma$ and $\tau$
    agree outside $D$ and $0$ otherwise, which is the element of $K\otimes1$ at
    $(g\sigma,g\tau)$. The tripartite coordinates satisfy these hypotheses for each of
    the four regions.
\end{proof}

\begin{lemma}[Decorrelation on a system of sites]
    \label{lem:commuting_parent_decorrelation}
    \lean{ManyBody.isDecorrelated_reindex_iff}
    \leanok
    \uses{def:sites_place, def:commuting_parent_coordinates}
    Following~\cite[Definition~D.1, lines 2187--2191]{Cirac2016MPDO_arXiv}, the regions
    $a$ and $b$ are decorrelated in an operator $P$ on
    $\mathcal H_a\otimes\mathcal H_m\otimes\mathcal H_b$ when
    $P(O_a\otimes1\otimes1)(1-P)(1\otimes1\otimes O_b)P=0$ for all complex matrices $O_a$
    on $\mathcal H_a$ and $O_b$ on $\mathcal H_b$. The source tests Hermitian $O_a,O_b$;
    since every matrix is $O=O'+iO''$ with $O',O''$ Hermitian, the two forms agree. For an
    operator $P$ on $\mathcal H_X$, the regions are decorrelated in $P$ read in the
    tripartite coordinates if and only if $P\,O_a\,(1-P)\,O_b\,P=0$ for all operators
    $O_a$ placed on $a$ and $O_b$ placed on $b$.
\end{lemma}
\begin{proof}\leanok
    \uses{lem:commuting_parent_placement, lem:sites_place_hom}
    Reading matrices in the tripartite coordinates is an isomorphism of algebras, and by
    Lemma~\ref{lem:commuting_parent_placement} it carries the placements on $a$ and on
    $b$ to $O_a\otimes1\otimes1$ and $1\otimes1\otimes O_b$.
\end{proof}

\begin{definition}[Commuting parent Hamiltonian]
    \label{def:commuting_parent_hamiltonian}
    \lean{ManyBody.HasCommutingParentHamiltonian}
    \leanok
    \uses{def:dynamics_interaction, def:hamiltonians_zero_space}
    An operator $P$ on $\mathcal H_X$ has a commuting parent Hamiltonian on
    $V\setminus b$ and $V\setminus a$ when there is an interaction $\Phi=(h_0,h_1)$ with
    $h_0$ supported on $V\setminus b$ and $h_1$ on $V\setminus a$, whose terms are
    commuting orthogonal projections, and whose zero space is $\operatorname{ran}P$.
    This is~\cite[Definition~D.2, lines 2205--2218]{Cirac2016MPDO_arXiv}: the terms are
    $1-P_{am}$ and $1-P_{mb}$ for the projections onto the two local ground spaces, and
    the zero space is the intersection of the two placed ground spaces.
\end{definition}

\begin{lemma}[The ground projection is the product]
    \label{lem:commuting_parent_product}
    \lean{ManyBody.HasCommutingParentHamiltonian.eq_mul}
    \leanok
    \uses{def:commuting_parent_hamiltonian}
    Let $P$ be an orthogonal projection, and let $\Phi=(h_0,h_1)$ be as in
    Definition~\ref{def:commuting_parent_hamiltonian}. Then
    $(1-h_0)(1-h_1)=P=(1-h_1)(1-h_0)$.
\end{lemma}
\begin{proof}\leanok
    \uses{def:hamiltonians_zero_space}
    The operators $1-h_0$ and $1-h_1$ are commuting orthogonal projections, so their
    product is an orthogonal projection whose range consists of the vectors fixed by
    both, $\operatorname{ran}(1-h_0)\cap\operatorname{ran}(1-h_1)=\ker h_0\cap\ker h_1$.
    This is the zero space of $\Phi$, which is $\operatorname{ran}P$. Orthogonal
    projections with the same range are equal, and the factors commute.
\end{proof}

\begin{theorem}[Commuting parent Hamiltonians and decorrelation]
    \label{thm:commuting_parent_iff_decorrelated}
    \lean{ManyBody.HasCommutingParentHamiltonian.isDecorrelated,
        ManyBody.hasCommutingParentHamiltonian_of_isDecorrelated,
        ManyBody.hasCommutingParentHamiltonian_iff_isDecorrelated}
    \leanok
    \uses{def:commuting_parent_hamiltonian, lem:commuting_parent_decorrelation}
    For disjoint regions $a$ and $b$, an orthogonal projection $P$ on $\mathcal H_X$ has
    a commuting parent Hamiltonian on $V\setminus b$ and $V\setminus a$ if and only if
    $a$ and $b$ are decorrelated in $P$, read in the tripartite coordinates.
    This is~\cite[Proposition~D.3, lines 2221--2289]{Cirac2016MPDO_arXiv}.
\end{theorem}
\begin{proof}\leanok
    \uses{lem:commuting_parent_product, lem:commuting_parent_placement,
        thm:sites_place_range, lem:sites_place_hom}
    Read every operator in the tripartite coordinates. Suppose $P$ has a commuting
    parent Hamiltonian $(h_0,h_1)$. By Theorem~\ref{thm:sites_place_range} and
    Lemma~\ref{lem:commuting_parent_placement}, $S_{am}=1-h_0$ is $G\otimes1$ and
    $S_{mb}=1-h_1$ is $1\otimes G'$, so $S_{mb}$ commutes with $O_a\otimes1\otimes1$ and
    $S_{am}$ with $1\otimes1\otimes O_b$. By Lemma~\ref{lem:commuting_parent_product},
    $P=S_{am}S_{mb}=S_{mb}S_{am}$, so $S_{mb}PS_{am}=(S_{mb}S_{am})^2=P^2=P$ and
    $S_{mb}(1-P)S_{am}=P-P=0$. Hence
    \begin{align*}
        P(O_a\otimes1\otimes1)(1-P)(1\otimes1\otimes O_b)P
        =S_{am}(O_a\otimes1\otimes1)\bigl[S_{mb}(1-P)S_{am}\bigr]
        (1\otimes1\otimes O_b)S_{mb}=0.
    \end{align*}
    Conversely, suppose $a$ and $b$ are decorrelated in $P$. Let $G_{am}$ and $G_{mb}$ be
    the projections onto the supports of the marginals of $P$ on $V\setminus b$ and on
    $V\setminus a$. By~\cite[lines 2225--2277]{Cirac2016MPDO_arXiv}, decorrelation gives
    $(1\otimes G_{mb})(G_{am}\otimes1)=P=(G_{am}\otimes1)(1\otimes G_{mb})$. Placing
    $G_{am}$ on $V\setminus b$ and $G_{mb}$ on $V\setminus a$ gives orthogonal
    projections that commute and whose product is $P$, and the interaction with terms
    $h_0=1-G_{am}$ and $h_1=1-G_{mb}$, placed, has commuting orthogonal projections as
    terms. A vector is annihilated by both terms exactly when it is fixed by both placed
    support projections, hence by their product $P$; conversely, a vector fixed by $P$
    is fixed by each factor, since each factor is idempotent and the factors commute.
    Thus the zero space is $\operatorname{ran}P$.
\end{proof}
